%
% reversi_240_feature_di_base.tex
%
% This file is part of the reversi program
% http://github.com/rcrr/reversi-v2
%
% Author Roberto Corradini mailto:rob_corradini@yahoo.it
% Copyright 2026 Roberto Corradini. All rights reserved.
%
% License
%
% This program is free software; you can redistribute it and/or modify it
% under the terms of the GNU General Public License as published by the
% Free Software Foundation; either version 3, or (at your option) any
% later version.
%
% This program is distributed in the hope that it will be useful,
% but WITHOUT ANY WARRANTY; without even the implied warranty of
% MERCHANTABILITY or FITNESS FOR A PARTICULAR PURPOSE. See the
% GNU General Public License for more details.
%
% You should have received a copy of the GNU General Public License
% along with this program; if not, write to the Free Software
% Foundation, Inc., 59 Temple Place - Suite 330, Boston, MA  02111-1307, USA
% or visit the site <http://www.gnu.org/licenses/>.
%

%-----------------------------------------------------------------------------
% Capitolo: Le feature di base.
%
% Prima campagna sperimentale del modello logistico: i quattro modelli
% minimali (l'intercetta; l'intercetta + LMC; l'intercetta + ALMC; la
% coppia completa), addestrati e validati sullo stesso corpus ec20, e
% la prima lettura comparata dei risultati (il contributo marginale
% dell'ALMC, l'assenza di overfitting).
%-----------------------------------------------------------------------------

\part{Esperienze con il modello logistico}

\chapter{Le feature di base}
\label{ch:feature-base}

Fino a qui il libro ha descritto \emph{come} il modello logistico viene
costruito: la teoria, la pipeline, i suoi livelli. Questa parte inverte
la prospettiva: impieghiamo il modello per caratterizzare il contributo
delle diverse feature disponibili, e di seguito assemblare modelli via
via più capaci.

Il metodo è quello della campagna sperimentale che questa parte
documenta: si costruiscono modelli \emph{minimali}, ciascuno con una
sola famiglia di feature in più rispetto al precedente, e si
confrontano i risultati di validazione. Tutti e quattro i modelli sono
addestrati e validati sullo \emph{stesso} corpus, con gli stessi
iparametri (frequency cut-off 1, logit clipping 0.05, regolarizzazione
ridge 0.01), in modo che ogni differenza di merito sia attribuibile
alla sola feature set.

\section{Il corpus e i quattro modelli}

Il corpus è il set di posizioni REGAB alla finestra ec20 (\emph{empty
count} 20): 199\,932 posizioni di addestramento e 20\,000 di
validazione, estratte dal database con i soli filtri di batch e di
stato. Le due parti sono \emph{disgiunte}: l'incrocio conta zero
posizioni comuni (0.00\%), quindi la validazione misura la capacità di
generalizzazione, non la memoria.

La tabella~\ref{tab:feature-base-modelli} elenca i quattro modelli. La
progressione è a scalini: si parte dal modello debole per eccellenza
(l'intercetta da sola) e si aggiunge una famiglia di feature per
volta.

\begin{table}[ht]
\centering
\begin{tabular}{llcl}
\toprule
ID & nome & feature set & parametri ($W$) \\
\midrule
0000 & baseline & intercetta & 1 \\
0001 & +\,LMC & intercetta + LMC & 22 \\
0002 & +\,ALMC & intercetta + ALMC & 22 \\
0003 & +\,LMC\,+\,ALMC & intercetta + LMC + ALMC & 43 \\
\bottomrule
\end{tabular}
\caption{I quattro modelli della prima campagna sperimentale (corpus ec20).}
\label{tab:feature-base-modelli}
\end{table}

LMC e ALMC sono le due mobility standard a mask piena descritte nel
capitolo~\ref{ch:feature}: il numero di mosse legali di chi muove, e il
numero di mosse che \emph{avrebbe} l'avversario se toccasse a lui.
Ogni mobility a mask piena ammette 65 configurazioni possibili (il
conto corre da 0 a 64); con la soglia di frequenza in uso, solo le
configurazioni osservate nel corpus restano nel modello, e il conteggio
effettivo è 21 per mobility. Da qui i 22 parametri dei modelli 0001 e
0002 (l'intercetta più 21 pesi) e i 43 del modello 0003 (l'intercetta
più 21 più 21).

\section{I risultati}

La tabella~\ref{tab:feature-base-risultati} raccoglie i risultati di
validazione dei quattro modelli.

\begin{table}[ht]
\centering
\begin{tabular}{lrrrr}
\toprule
metrica & 0000 & 0001 & 0002 & 0003 \\
\midrule
RMSE (y)           & 26.28  & 21.81  & 22.47  & \textbf{20.13} \\
MAE (y)            & 21.63  & 17.90  & 18.43  & \textbf{16.47} \\
Loss ($\mathrm{MSE}/2$) & $1.707\times10^{-2}$ & $1.176\times10^{-2}$ & $1.248\times10^{-2}$ & \textbf{$1.002\times10^{-2}$} \\
parametri $W$      & 1      & 22     & 22     & 43 \\
parametri dense    & 1      & 66     & 66     & 131 \\
iterazioni L-BFGS  & 5      & 40     & 42     & 60 \\
gap di generalizzazione & $-0.13$ & $-0.20$ & $-0.11$ & $-0.20$ \\
\bottomrule
\end{tabular}
\caption{Risultati di validazione dei quattro modelli (corpus ec20, 20\,000
posizioni). Il valore in grassetto è il migliore di ogni riga.}
\label{tab:feature-base-risultati}
\end{table}

La cosa che salta agli occhi per prima è la prima colonna: il modello
0000, privo di ogni feature, riproduce \emph{per costruzione} la
deviazione standard della popolazione di validazione (26.28 contro
26.28): predice la media, e il suo errore \emph{è} la dispersione dei
valori di gioco. È il riferimento contro cui misurare tutto il resto.

Rispetto a quel riferimento, i guadagni in RMSE sono:

\begin{itemize}
  \item LMC da sola (0001): $-17.0\%$;
  \item ALMC da sola (0002): $-14.4\%$, poco più debole;
  \item la coppia completa (0003): $-23.6\%$.
\end{itemize}

\section{Contributo marginale e complementarità}

La domanda interessante non è quanto vale l'ALMC da sola, ma quanto
vale \emph{in presenza} della LMC. La risposta si legge sottraendo:
la coppia porta l'errore da 21.81 a 20.13, un contributo marginale
netto di $-1.68$ in RMSE ($-7.7\%$ sul modello 0001) al prezzo di 21
parametri in più.

Il contributo marginale è più piccolo del contributo da sola ($-1.68$
contro $-3.81$): l'informazione di ALMC sovrappone in parte quella di
LMC — è plausibile, le due quantità contano mosse sulla stessa board —
ma una frazione sostanziale è informazione nuova, quante mosse
\emph{perderebbe} l'avversario. Le due mobility sono
quindi \emph{complementari}, non ridondanti: entrambe restano nel
feature set.

Due controlli di sanità, validi per tutti e quattro i modelli. Il gap
di generalizzazione (RMSE di training meno RMSE di validazione) resta
\emph{negativo}: la validazione batte sempre l'addestramento, e non c'è
alcun segnale di overfitting fino a 43 parametri. E il controllo di
client — il round-trip tra il file dei pesi e le predizioni in
memoria — è esatto in tutti e quattro i run (differenza massima
$0.0\times10^{0}$).
